%! Inverse Functions — Unit 1, Lesson 1.6 — Exit Ticket (ANSWER KEY)
\documentclass[10pt]{article}
\usepackage{precalculus-article}
\usepackage{precalculus-key}
\newcommand{\TallMath}[1]{$\displaystyle #1\rule[-1.4em]{0pt}{3.2em}$}

\begin{document}

\pageheader{Unit 1, Lesson 1.6}{Exit Ticket --- Answer Key}

\vspace{0.12in}

\begin{enumerate}[itemsep=12pt]

\item A function $f$ satisfies $f(6) = -2$.

\begin{enumerate}[label=(\alph*), itemsep=5pt]
  \item $f^{-1}(-2) = $ \ans{6}
  \item In one sentence, say why $f^{-1}(-2)$ is not $-\dfrac{1}{2}$.

        \vspace{0.05in}
        \ansline{$f^{-1}$ gives back the question that produced $-2$, which was 6; the reciprocal is just arithmetic done to the answer.}
\end{enumerate}

\item Let $f(x) = 2x - 5$. Find $f^{-1}(x)$ by swapping and solving.

\begin{work}
              y &= 2x - 5 \\
              x &= 2y - 5 \\
          x + 5 &= 2y \\
  \frac{x+5}{2} &= y
\end{work}

$f^{-1}(x) = $ \ans{$\dfrac{x+5}{2}$}

\smallskip
Now use it: $f^{-1}(9) = $ \ans{7}, and checking forwards,
$f($ \ans{7} $) = 9$.

\item The table gives every value of a function $h$.

\smallskip
\renewcommand{\arraystretch}{1.4}
\begin{center}
\begin{tabular}{|c|c|c|c|c|}
\hline
$x$ & 1 & 2 & 3 & 4 \\
\hline
$h(x)$ & 5 & 8 & 5 & 12 \\
\hline
\end{tabular}
\end{center}

\smallskip
$h$ has no inverse. The answer \ans{5} came back from two different
questions, $x = $ \ans{1} and $x = $ \ans{3}, so $h$ is not
\ans{one-to-one}.

\end{enumerate}

\end{document}
